% GENERATED FILE. Edit canonical lessons, then run npm run generate:books.
% canonical-source-hash: fnv1a64:e3de5f0eabea58f6
% canonical-lessons: KA-C164-aplod-madu, KA-C164-klik-madu, KA-C164-alisu, KA-C164-an-madu, KA-C164-aph-madu

\chapter{To upload, To click, To delete, To switch on, To switch off}
\label{ch:ka-to-upload-to-click-to-delete-to-switch-on-to-switch-off}
\hlchaptermodality{164}

\begin{chapteropening}
\textbf{By the end of this chapter:} \emph{I can say to upload, to click, to delete, to switch on and to switch off.}

Complete the last lesson of chapter 164: I can say to upload, to click, to delete, to switch on and to switch off.
\end{chapteropening}

\section[\texorpdfstring{\kn{ಅಪ್ಲೋಡ್} \kn{ಮಾಡು} (aplōḍ māḍu)}{ಅಪ್ಲೋಡ್ ಮಾಡು (aplōḍ māḍu)}]{\kn{ಅಪ್ಲೋಡ್} \kn{ಮಾಡು} (aplōḍ māḍu) — to upload}

\label{lesson:KA-C164-aplod-madu}



\begin{quote}
Before the new one: say the Kannada for to get along, then the Kannada for to plant.
\end{quote}

\subsection*{You'll want to know: \kn{ಅಪ್ಲೋಡ್} \kn{ಮಾಡು}}
\textbf{\kn{ಅಪ್ಲೋಡ್} \kn{ಮಾಡು}} — \emph{aplōḍ māḍu} — \textquotedblleft{}to upload\textquotedblright{}.

\begin{grammarlens}[title={What you've built}]
One more word for this chapter.
\end{grammarlens}

\subsection*{Your turn}
\begin{itemize}
\raggedright
  \item \emph{Say it:} \emph{aplōḍ māḍu}
  \item \emph{Say it:} \emph{aplōḍ māḍu}, once more
  \item \emph{Say it:} say \emph{aplōḍ māḍu}
  \item \emph{Recall:} say the Kannada for to get along, then the Kannada for to plant, then say \emph{aplōḍ māḍu} again
\end{itemize}

\begin{tcolorbox}[breakable,colback=teal!4,colframe=teal!35!black,title={Before you move on}]
What does \kn{ಅಪ್ಲೋಡ್} \kn{ಮಾಡು} mean? (\textquotedblleft{}To upload\textquotedblright{}.) Say it once more.
\end{tcolorbox}
\section[\texorpdfstring{\kn{ಕ್ಲಿಕ್} \kn{ಮಾಡು} (klik māḍu)}{ಕ್ಲಿಕ್ ಮಾಡು (klik māḍu)}]{\kn{ಕ್ಲಿಕ್} \kn{ಮಾಡು} (klik māḍu) — to click}

\label{lesson:KA-C164-klik-madu}



\begin{quote}
Before the new one: say the Kannada for to plant, then the Kannada for to upload.
\end{quote}

\subsection*{You'll want to know: \kn{ಕ್ಲಿಕ್} \kn{ಮಾಡು}}
\textbf{\kn{ಕ್ಲಿಕ್} \kn{ಮಾಡು}} — \emph{klik māḍu} — \textquotedblleft{}to click\textquotedblright{}.

\begin{grammarlens}[title={What you've built}]
One more word for this chapter.
\end{grammarlens}

\subsection*{Your turn}
\begin{itemize}
\raggedright
  \item \emph{Say it:} \emph{klik māḍu}
  \item \emph{Say it:} \emph{klik māḍu}, once more
  \item \emph{Say it:} say \emph{klik māḍu}
  \item \emph{Recall:} say the Kannada for to plant, then the Kannada for to upload, then say \emph{klik māḍu} again
\end{itemize}

\begin{tcolorbox}[breakable,colback=teal!4,colframe=teal!35!black,title={Before you move on}]
What does \kn{ಕ್ಲಿಕ್} \kn{ಮಾಡು} mean? (\textquotedblleft{}To click\textquotedblright{}.) Say it once more.
\end{tcolorbox}
\section[\texorpdfstring{\kn{ಅಳಿಸು} (aḷisu)}{ಅಳಿಸು (aḷisu)}]{\kn{ಅಳಿಸು} (aḷisu) — to delete, to erase}

\label{lesson:KA-C164-alisu}



\begin{quote}
Before the new one: say the Kannada for to upload, then the Kannada for to click.
\end{quote}

\subsection*{You'll want to know: \kn{ಅಳಿಸು}}
\textbf{\kn{ಅಳಿಸು}} — \emph{aḷisu} — \textquotedblleft{}to delete, to erase\textquotedblright{}.

\begin{grammarlens}[title={What you've built}]
One more word for this chapter.
\end{grammarlens}

\subsection*{Your turn}
\begin{itemize}
\raggedright
  \item \emph{Say it:} \emph{aḷisu}
  \item \emph{Say it:} \emph{aḷisu}, once more
  \item \emph{Say it:} say \emph{aḷisu}
  \item \emph{Recall:} say the Kannada for to upload, then the Kannada for to click, then say \emph{aḷisu} again
\end{itemize}

\begin{tcolorbox}[breakable,colback=teal!4,colframe=teal!35!black,title={Before you move on}]
What does \kn{ಅಳಿಸು} mean? (\textquotedblleft{}To delete, to erase\textquotedblright{}.) Say it once more.
\end{tcolorbox}
\section[\texorpdfstring{\kn{ಆನ್} \kn{ಮಾಡು} (ān māḍu)}{ಆನ್ ಮಾಡು (ān māḍu)}]{\kn{ಆನ್} \kn{ಮಾಡು} (ān māḍu) — to switch on}

\label{lesson:KA-C164-an-madu}



\begin{quote}
Before the new one: say the Kannada for to click, then the Kannada for to delete.
\end{quote}

\subsection*{You'll want to know: \kn{ಆನ್} \kn{ಮಾಡು}}
\textbf{\kn{ಆನ್} \kn{ಮಾಡು}} — \emph{ān māḍu} — \textquotedblleft{}to switch on\textquotedblright{}.

\begin{grammarlens}[title={What you've built}]
One more word for this chapter.
\end{grammarlens}

\subsection*{Your turn}
\begin{itemize}
\raggedright
  \item \emph{Say it:} \emph{ān māḍu}
  \item \emph{Say it:} \emph{ān māḍu}, once more
  \item \emph{Say it:} say \emph{ān māḍu}
  \item \emph{Recall:} say the Kannada for to click, then the Kannada for to delete, then say \emph{ān māḍu} again
\end{itemize}

\begin{tcolorbox}[breakable,colback=teal!4,colframe=teal!35!black,title={Before you move on}]
What does \kn{ಆನ್} \kn{ಮಾಡು} mean? (\textquotedblleft{}To switch on\textquotedblright{}.) Say it once more.
\end{tcolorbox}
\section[\texorpdfstring{\kn{ಆಫ್} \kn{ಮಾಡು} (āph māḍu)}{ಆಫ್ ಮಾಡು (āph māḍu)}]{\kn{ಆಫ್} \kn{ಮಾಡು} (āph māḍu) — to switch off}

\label{lesson:KA-C164-aph-madu}



\begin{quote}
Before the new one: say the Kannada for to delete, then the Kannada for to switch on.
\end{quote}

\subsection*{You'll want to know: \kn{ಆಫ್} \kn{ಮಾಡು}}
\textbf{\kn{ಆಫ್} \kn{ಮಾಡು}} — \emph{āph māḍu} — \textquotedblleft{}to switch off\textquotedblright{}.

\begin{grammarlens}[title={What you've built}]
One more word for this chapter.
\end{grammarlens}

\subsection*{Your turn}
\begin{itemize}
\raggedright
  \item \emph{Say it:} \emph{āph māḍu}
  \item \emph{Say it:} \emph{āph māḍu}, once more
  \item \emph{Say it:} say \emph{āph māḍu}
  \item \emph{Recall:} say the Kannada for to delete, then the Kannada for to switch on, then say \emph{āph māḍu} again
\end{itemize}

\begin{tcolorbox}[breakable,colback=teal!4,colframe=teal!35!black,title={Before you move on}]
What does \kn{ಆಫ್} \kn{ಮಾಡು} mean? (\textquotedblleft{}To switch off\textquotedblright{}.) Say it once more.
\end{tcolorbox}
